% Exam Template for UMTYMP and Math Department courses
%
% Using Philip Hirschhorn's exam.cls: http://www-math.mit.edu/~psh/#ExamCls
%
% run pdflatex on a finished exam at least three times to do the grading table on front page.
%
%%%%%%%%%%%%%%%%%%%%%%%%%%%%%%%%%%%%%%%%%%%%%%%%%%%%%%%%%%%%%%%%%%%%%%%%%%%%%%%%%%%%%%%%%%%%%%

% These lines can probably stay unchanged, although you can remove the last
% two packages if you're not making pictures with tikz.
\documentclass[11pt]{exam}
\RequirePackage{amssymb, amsfonts, amsmath, latexsym, verbatim, xspace, setspace}
\RequirePackage{tikz, pgflibraryplotmarks}

% By default LaTeX uses large margins.  This doesn't work well on exams; problems
% end up in the "middle" of the page, reducing the amount of space for students
% to work on them.
\usepackage[margin=1in]{geometry}
\usepackage{enumerate}
\usepackage{amsthm}

\theoremstyle{definition}
\newtheorem{soln}{Solution}

% Here's where you edit the Class, Exam, Date, etc.
\newcommand{\class}{Math 350 Section 1}
\newcommand{\term}{Fall 2024}
\newcommand{\examnum}{Exam III}
\newcommand{\examdate}{November 19, 2024}
\newcommand{\timelimit}{1 Hour 50 Minutes}
\newcommand{\ol}[1]{\overline{#1}}

% For an exam, single spacing is most appropriate
\singlespacing
% \onehalfspacing
% \doublespacing

% For an exam, we generally want to turn off paragraph indentation
\parindent 0ex

\begin{document} 

% These commands set up the running header on the top of the exam pages
\pagestyle{head}
\firstpageheader{}{}{}
\runningheader{\class}{\examnum\ - Page \thepage\ of \numpages}{\examdate}
\runningheadrule

\begin{flushright}
\begin{tabular}{p{2.8in} r l}
\textbf{\class} & \textbf{Name (Print):} & \makebox[2in]{\hrulefill}\\
\textbf{\term} &&\\
\textbf{\examnum} & \textbf{Student ID:}&\makebox[2in]{\hrulefill}\\
\textbf{\examdate} &&\\
\textbf{Time Limit: \timelimit} % & Teaching Assistant & \makebox[2in]{\hrulefill}
\end{tabular}\\
\end{flushright}
\rule[1ex]{\textwidth}{.1pt}


This exam contains \numpages\ pages (including this cover page) and
\numquestions\ problems.  Check to see if any pages are missing.  Enter
all requested information on the top of this page, and put your initials
on the top of every page, in case the pages become separated.\\

You may \textit{not} use your books or notes on this exam.

You are required to show your work on each problem on this exam.  The following rules apply:\\

\begin{minipage}[t]{3.7in}
\vspace{0pt}
\begin{itemize}

%\item \textbf{If you use a ``fundamental theorem'' you must indicate this} and explain
%why the theorem may be applied.

\item \textbf{Organize your work}, in a reasonably neat and coherent way, in
the space provided. Work scattered all over the page without a clear ordering will 
receive very little credit.  

\item \textbf{Mysterious or unsupported answers will not receive full
credit}.  A correct answer, unsupported by calculations, explanation,
or algebraic work will receive no credit; an incorrect answer supported
by substantially correct calculations and explanations might still receive
partial credit.  This especially applies to limit calculations.

\item If you need more space, use the back of the pages; clearly indicate when you have done this.

\item \textbf{Box Your Answer} where appropriate, in order to clearly indicate what you consider the answer to the question to be.

\item \textbf{Default metric}: unless otherwise stated, the set $\mathbb{R}$ or $\mathbb{R}^n$ will be considered a metric space with the Euclidean metric
\end{itemize}

Do not write in the table to the right.
\end{minipage}
\hfill
\begin{minipage}[t]{2.3in}
\vspace{0pt}
%\cellwidth{3em}
\gradetablestretch{2}
\vqword{Problem}
\addpoints % required here by exam.cls, even though questions haven't started yet.	
\gradetable[v]%[pages]  % Use [pages] to have grading table by page instead of question

\end{minipage}
\newpage % End of cover page

%%%%%%%%%%%%%%%%%%%%%%%%%%%%%%%%%%%%%%%%%%%%%%%%%%%%%%%%%%%%%%%%%%%%%%%%%%%%%%%%%%%%%
%
% See http://www-math.mit.edu/~psh/#ExamCls for full documentation, but the questions
% below give an idea of how to write questions [with parts] and have the points
% tracked automatically on the cover page.
%

%%%%%%%%%%%%%%%%%%%%%%%%%%%%%%%%%%%%%%%%%%%%%%%%%%%%%%%%%%%%%%%%%%%%%%%%%%%%%%%%%%%%%

\begin{questions}

\addpoints

\question[10]\mbox{}
\textbf{TRUE or FALSE!}  Write  TRUE if the statement is true.  Otherwise, write FALSE.  Your response should be in ALL CAPS.  No justification is required.
\begin{enumerate}[(a)]
\item there is a continuous, real-valued function on $[0,1]$ whose range is $[0,\infty)$
\vspace{1.2in}
\item a sequence $\{x_n\}$ of numbers in $\mathbb{R}$ which is monotone increasing and bounded above converges
\vspace{1.2in}
\item if $f: (a,b)\rightarrow\mathbb{R}$ is differentiable at a point $c\in (a,b)$, then it is also continuous at $c$
\vspace{1.2in}
\item a Cauchy sequence in $\mathbb{R}$, but with the discrete metric, will converge
\vspace{1.2in}
\item the equation $x^{25}+2x+2=0$ has a solution in the interval $[-1,0]$
\vspace{1.2in}
\end{enumerate}

\newpage
\question[10]\mbox{}

\begin{enumerate}[(a)]
\item
Let $(S,d_S)$ and $(T,d_T)$ be metric spaces and $A\subseteq S$.
Write down what it means for a function $f: A\rightarrow T$ to be continuous at a point $a\in A$.
\vspace{2in}
\item 
If $f: A\rightarrow T$ is a function, write down the definition of the preimage $f^{-1}(V)$ of a subset $V\subseteq T$.
\vspace{2in}
\item 
Prove that if $f: A\rightarrow T$ is continuous at every point of $A$, then for any open subsets $V\subseteq T$, the preimage $f^{-1}(V)$ is open.
\end{enumerate}

\newpage
\question[10]\mbox{}

\begin{enumerate}[(a)]
\item  Let $(S,d_S)$ and $(T,d_T)$ be metric spaces and $A\subseteq S$.  If $a\in S$ is accumulation point of $A$ and $f: A\rightarrow T$ is a function, write the definition of $\lim_{x\rightarrow a} f(x) = L$.
\vspace{1in}
\item  Prove that if  $\lim_{x\rightarrow a} f(x) = L$, then for any sequence $\{x_n\}$ of elements in $A$ converging to $a$ we have $\lim_{n\rightarrow\infty} f(x_n)=L$
\vspace{3in}
\item  Consider the function $f:\mathbb{R}\rightarrow\mathbb{R}$ defined by
$$f(x) = \left\lbrace\begin{array}{cc}
1, & x\in\mathbb{Q}\\
0, & x\notin\mathbb{Q}
\end{array}\right.$$
Does the limit $\lim_{x\rightarrow 0} f(x)$ exist?  Explain.
\end{enumerate}

\newpage
\question[10]\mbox{}
\begin{enumerate}[(a)]
\item Write down the definition of a metric space $(M,d)$ being disconnected, as well as the definition of a subset $A\subseteq M$ being disconnected.
\vspace{1in}
\item Prove that if $A,B\subseteq M$ are connected and overlap, then $A\cup B$ is also connected.
\vspace{3in}
\item Give an example showing that if $A,B\subseteq M$ are connected and overlap, then $A\cap B$ may not be connected.
\end{enumerate}

\newpage

\newpage
\question[10]\mbox{}
\begin{enumerate}[(a)]
\item Let $f$ be a real-valued function on $(a,b)$.  Write down what it means for $f$ to be differentiable at $c\in (a,b)$.
\vspace{2in}
\item Carefully state the Helper Theorem we used in class.
\vspace{2in}
\item Use the Helper Theorem to prove the Product Rule for derivatives.
\end{enumerate}

\newpage
\question[10]\mbox{}
Let $f$ and $\alpha$ be a real-valued function on an interval $[a,b]$.
\begin{enumerate}[(a)]
\item  Given a partition $P$ of $[a,b]$ and sample points $\{t_k\}$, define $S(P,f,\alpha,\{t_k\})$
\vspace{1in}
\item  Write the definition of the Riemann integral of $f$ on $[a,b]$ with respect to $\alpha$.
\vspace{2in}
\item  Let $\alpha(x)=x$ and $f$ be the function
$$f(x) = \left\lbrace\begin{array}{cc}
0, & x < 1\\
1, & 1 \leq x < 2\\
2, & 2\leq x
\end{array}\right.$$
Use the definition to show that $f$ is Riemann integrable with respect to $\alpha$ on $[0,3]$ and calculate $\int_0^3fd\alpha$
\end{enumerate}

\newpage
\question[10]\mbox{}
\begin{enumerate}[(a)]
\item Use the definition of a derivative to prove that $f(x) = \sqrt{x}$ is differentiable and satisfies $f'(a) = \frac{1}{2\sqrt{a}}$ for all $a > 0$
\vspace{3in}
\item Prove that for any $a>0$
$$\frac{1}{2\sqrt{a+1}} < \sqrt{a+1}-\sqrt{a} < \frac{1}{2\sqrt{a}}$$
Hint: think about the Mean Value Theorem for $f(x) = \sqrt{x}$
\end{enumerate}


\end{questions}

\end{document}

